% Three genuine midpoint-chord examples on the same hyperbola and scale.
\begin{tikzpicture}[x=.74cm,y=.74cm,font=\small,>=stealth]
\node[font=\small\bfseries] at (1,2.9) {三种“确实存在”的情形：统一取 $x^2/4-y^2=1$};
\draw[CStatement,thick,domain=-1.6:1.6,samples=100] plot ({2*sqrt(1+\x*\x)},\x);
\draw[CStatement,thick,domain=-1.6:1.6,samples=100] plot ({-2*sqrt(1+\x*\x)},\x);
\draw[gray,->] (-4.05,0)--(4.1,0);
\draw[gray,->] (0,-1.9)--(0,1.95);
\draw[teal,very thick] ({-2*sqrt(2)},1)--({2*sqrt(2)},1);
\fill ({-2*sqrt(2)},1) circle(2pt) node[below left] {$A$};
\fill ({2*sqrt(2)},1) circle(2pt) node[below right] {$B$};
\fill[orange!80!black] (0,1) circle(2.2pt);
\node[orange!80!black] at (.8,1.6) {$M$};
\node[anchor=west,align=left,text width=5cm] at (4.8,.2) {
\textbf{情形一：$S<0$：异支弦}\\[4pt]
$M=(0,1),\quad S=-1<0$。\\[4pt]
$A=(-2\sqrt2,1)$，$B=(2\sqrt2,1)$。\\[4pt]
端点在不同支，$AM=MB$。
};
\begin{scope}[yshift=-4.5cm]
\draw[CStatement,thick,domain=-1.6:1.6,samples=100] plot ({2*sqrt(1+\x*\x)},\x);
\draw[CStatement,thick,domain=-1.6:1.6,samples=100] plot ({-2*sqrt(1+\x*\x)},\x);
\draw[gray,->] (-4.05,0)--(4.1,0);
\draw[gray,->] (0,-1.9)--(0,1.95);
\draw[teal,very thick] (3,{-sqrt(5)/2})--(3,{sqrt(5)/2});
\fill (3,{sqrt(5)/2}) circle(2pt) node[above left] {$A$};
\fill (3,{-sqrt(5)/2}) circle(2pt) node[below left] {$B$};
\fill[orange!80!black] (3,0) circle(2.2pt);
\node[orange!80!black] at (3.65,.5) {$M$};
\node[anchor=west,align=left,text width=5cm] at (4.8,.2) {
\textbf{情形二：$S>1$：同支弦}\\[4pt]
$M=(3,0),\quad S=9/4>1$。\\[4pt]
$A=(3,\sqrt5/2)$，\\[4pt]
$B=(3,-\sqrt5/2)$。\\[4pt]
端点同支，$AM=MB$；本例为竖直弦。
};
\end{scope}
\begin{scope}[yshift=-9cm]
\draw[CStatement,thick,domain=-1.6:1.6,samples=100] plot ({2*sqrt(1+\x*\x)},\x);
\draw[CStatement,thick,domain=-1.6:1.6,samples=100] plot ({-2*sqrt(1+\x*\x)},\x);
\draw[gray,->] (-4.05,0)--(4.1,0);
\draw[gray,->] (0,-1.9)--(0,1.95);
\draw[teal,very thick] (-2,0)--(2,0);
\draw[CProof,dashed,thick] ({-4/sqrt(3)},{-1/sqrt(3)})--({4/sqrt(3)},{1/sqrt(3)});
\fill (-2,0) circle(2pt) node[above left] {$A$};
\fill (2,0) circle(2pt) node[below right] {$B$};
\fill[CProof] ({-4/sqrt(3)},{-1/sqrt(3)}) circle(1.7pt);
\fill[CProof] ({4/sqrt(3)},{1/sqrt(3)}) circle(1.7pt);
\fill[orange!80!black] (0,0) circle(2.2pt);
\node[orange!80!black] at (1.4,-1.25) {$M=O$};
\node[anchor=west,align=left,text width=5cm] at (4.8,.2) {
\textbf{情形三：$M=O$：中心例外}\\[4pt]
$S=0$，仍能有中心弦。\\[4pt]
实线 $y=0$，虚线 $y=x/4$，\\[4pt]
中点均为 $O$，方向不唯一。\\[4pt]
须有 $|k|<1/2$，并非任意方向。
};
\end{scope}
\end{tikzpicture}
